\documentclass{article}
\usepackage[left=2cm, right=2cm, top=2cm, bottom=2cm]{geometry}
\usepackage{graphicx}
\usepackage{amsmath}
\usepackage{amssymb}
\usepackage{amsthm}
\usepackage{fancyhdr}
\usepackage{verbatim}
\usepackage{listings}
\usepackage{xcolor}
\usepackage{pdfpages}
\usepackage{cancel}
\usepackage{physics}
\usepackage{esint}

\lstset{
    language=C++,
    basicstyle=\ttfamily\footnotesize,
    keywordstyle=\color{blue},
    commentstyle=\color{green},
    stringstyle=\color{red},
    numbers=left,
    numberstyle=\tiny\color{gray},
    frame=single,
    breaklines=true,
    captionpos=b,
}


\title{PHYS 487: Lecture \# 17}
\author{Cliff Sun}

\newtheorem{theorem}{Theorem}[section]
\newtheorem{lemma}[theorem]{Lemma}
\newtheorem{definition}[theorem]{Definition}
\newtheorem{conjecture}[theorem]{Conjecture}
\newtheorem{proposition}[theorem]{Proposition}
\newtheorem{corollary}[theorem]{Corollary}
\newtheorem{one minute paper}[theorem]{One Minute Paper}

\pagestyle{fancy}
\lhead{\textbf{\thepage}\ \ \nouppercase{\rightmark}}
\chead{PHYS 487: Lecture \# 17}
\rhead{Cliff Sun}

\begin{document}

\maketitle

\section*{Lecture Span}
\begin{enumerate}
    \item Spontaneous Emission and Fermi's Golden rule
\end{enumerate}

\section*{Einstein A \& B Coefficients}

Consider a box with a bunch of two level atoms. Ground state $\ket{a}$, excited state $\ket{b}$. Then, the total number of atoms is 

\begin{equation}
    N_{\Sigma} = N_A + N_B
\end{equation}

Where $N_i$ is the number of atoms in the i-th state. Assume that atoms can decay spontaneously with rate $A$ (atoms decaying per unit time is then $AN_b(t)$). 
Stimulated emission has a rate of $B_{b,a}\rho(\omega_0)$ ($\rho$ is the energy density per frequency). Then the number of atoms decaying due to stimulated emission is $B_{ba}N_b(t)\rho(\omega_0)$.
Absorption: $B_{ab}\rho(\omega_0)$, then number of excited atoms per time is $B_{ab}\rho(\omega_0)N_A(t)$. 

\par Then, we can create a rate equation

\begin{equation}
    \dot{N}_{b} = -N_B A - N_B B_{ba}\rho(\omega_0) + N_a B_{ab}\rho(\omega_0)
\end{equation}

We can make two assumptions. Firstly, the box's temperature is fixed, therefore $\rho(\omega_0)$ must be Boltzmann-distributed. Moreover, we can assume equilibrium, $\dot{N}_b = 0$. Then,
\begin{equation}
    \rho_{\omega_0} = \frac{A}{(N_A/N_B)B_{ab} - B_{ba}}
\end{equation}

Then from Boltzmann, we have that 

\begin{equation}
    \frac{N_a}{N_b} = \frac{\exp(-E_a/k_BT)}{\exp(-E_b/k_B T)} = \exp(\frac{\hbar \omega_0}{k_B T})
\end{equation}

Which then yields 

\begin{equation}
    \rho(\omega_0) = A\left( \exp\left( \frac{\hbar \omega_0}{k_B T}\right)B_{ab} - B_{ba} \right)^{-1}
\end{equation}

Planck's black body radiation formula states that 

\begin{equation}
    \rho(\omega) = \frac{\hbar}{\pi^2c^3}\frac{\omega^3}{e^{\hbar \omega/k_B T} - 1}
\end{equation}

You can find a solution if $B_{ab} = B_{ba}$, then 

\begin{equation}
    A = \frac{\omega_0^3\hbar}{\pi^2 c^3}B_{ba}
\end{equation}

With perturbation theory, 

\begin{equation}
    B_{ba} = \frac{\pi}{3\epsilon_0 \hbar^2}|p|^2, \quad A = \frac{\omega_0^3 |p|^2}{3\pi \epsilon_0 \hbar c^3}
\end{equation}

Consider a harmonic oscillator, starting on state $\ket{n}$. Then, the dipole moment is

\begin{equation}
    \hat{p} = \langle n | \hat{x} | n' \rangle
\end{equation}

And 

\begin{equation}
    \hat{x} = \sqrt{\frac{\hbar}{2m \omega}}\left( a_{+} + a_{-} \right)
\end{equation}

Also, assume $n' < n$. Then, we get a dipole moment 

\begin{equation}
    \hat{p} = q\sqrt{\frac{n\hbar}{2m \omega}}\delta_{n',n-1}
\end{equation}

Then, 

\begin{equation}
    A_n = \frac{nq^2 \omega^2}{6\pi \epsilon_0 mc^3} \equiv \frac{1}{\tau_n}
\end{equation}

Note, the radiated power is $$P = A_n \cdot \hbar \omega = \frac{q^2\omega^2}{6\pi \epsilon_0 mc^3}(E - \hbar \omega/2)$$. 

Classically, we have to use coherent states. (NOT STATIONARY STATES, since they are not classical)

\begin{equation}
    a_{-}\ket{\alpha} = \alpha\ket{\alpha}, \quad \ket{\alpha} = e^{-|\alpha|^2/2}\sum_{k}\frac{\alpha^k}{\sqrt{k!}}\ket{k}, \quad \langle a_{+}a_{-} \rangle = |\alpha|^2, \quad \langle H \rangle = \hbar \omega(|\alpha|^2 + 1/2)
\end{equation}

Probability in state $n$ is $$\langle n | \alpha \rangle^2 = e^{-|\alpha|^2/2}\frac{\alpha^{2n}}{n!}$$ 

Therefore, the probability of decaying (power) 

\begin{align*}
    P = \hbar \sum_n |\langle n | \alpha \rangle |^2 \cdot A_n &= \hbar \omega \sum_n | \langle n | \alpha \rangle |^2 \cdot n \cdot A_1 \\
    &= \hbar A_1 e^{-\langle \hat{n} \rangle}\sum_n \frac{\langle n \rangle^n}{n!}\cdot n \\
    &= \hbar \omega A_1 e^{-\langle n \rangle} \langle n \rangle \sum_{n}\frac{\langle n \rangle^{n-1}}{(n-1)!} \\
    &= \hbar \omega A_1 \langle n \rangle
\end{align*}

Note that $A_n = n \cdot A_1$. We can also compute the energy decay, note that 

\begin{equation}
    P = -\dot{E} = -\langle \dot{n} \rangle \hbar \omega = \hbar \omega A_1 \langle n \rangle = \langle n \rangle(t) = \langle\hat{n}\rangle(0)e^{-A_1t}
\end{equation}

\section*{Fermi's Golden Rule}

Given an initial state $\ket{i}$, then what's the rate that the initial state transitions to $\ket{f}$ for a single state: 

\begin{equation}
    P_{i\rightarrow f} = \frac{| \langle i | V | f \rangle|^2}{\hbar^2}\frac{\sin^2(\Delta\omega t/2)}{\Delta \omega^2}
\end{equation}

But what if $\ket{f}$ is a continuum of states? Then $P$ (probability of leaving the initial state) is

\begin{equation}
    P = \int dE_n p(E_n)\frac{| \langle i | V | n \rangle |^2}{\hbar^2}\frac{\sin^2(\Delta\omega t/2)}{\Delta \omega^2}
\end{equation}

For $t \rightarrow \infty$, then pull out $\rho(E_f)$. We get an approximation: 

\begin{equation}
    P \approx \frac{\pi | \langle i | V | f \rangle^2}{2\hbar}\rho(E_f)t
\end{equation}

Then, the rate is 

\begin{equation}
    R = \frac{\pi | \langle i | V | f \rangle^2}{2\hbar}\rho(E_f) = \frac{2\pi }{\hbar}|Vif/2|^2\rho(E_f)
\end{equation}

\end{document}